\scanpage{95}

\reversemarginpar
\margin{\begin{tikzpicture}[x=1cm,y=1cm,baseline=(current bounding box.north)]
\draw[line width=.25pt] (.04,0) .. controls (.00,-.20) and (.03,-1.1) .. (.10,-2.4) .. controls (.17,-3.6) and (.31,-5.2) .. (.28,-5.73);
\draw[line width=.25pt] (.28,-5.79) .. controls (.10,-5.81) and (-.05,-5.96) .. (-.24,-6.13);
\end{tikzpicture}}
\normalmarginpar

that $N(x)$ induces on $R^n$ the function $N_0(x)$, and that we have $N(\lambda x)=|\lambda|N(x)$ and $N(x)=0\Rightarrow x=0$, we need only prove $N(x+y)\leq N(x)+N(y)$ (\textcolor{red}{\uline{\textcolor{black}{is}}} valid even for $x,y\in c_0$.) For this, we \deleted{\unsure{regard} $(x_i)$ and $(y_i)$ as \unsure{the}}\inserted{observe that if} $u_x$ is the operation of multiplication by $x$ in $\ell^2$ (which is a compact operator for $x\in c_0$), we have \deleted{$u_{x+y}=u_x$} $\rho_i(u_x)=\rho_i(x)$, on the other hand \deleted{\unsure{that|$\rho_i$}} $u_{x+y}=u_x+u_y$, and the inequality to be proved then follows from number~4, theorem~4, corollary. At the same time, the \inserted{Schatten} norms \deleted{\unsure{ft}} thus associated with Schatten norms \deleted{\unsure{of}} $N_0$ on $R^n$ \deleted{give} demonstrate that there exists a \inserted{continuous} infinity of Schatten norms on $R^{(\underline N)}$ equivalent to the norm induced by $c_0$.

Let $N$ \deleted{\unsure{is}} be a Schatten norm on \deleted{$\unsure{R^I|R^n}$ \unsure{finite|finite}} $R^{(\underline N)}$. \deleted{Let $\unsure{\ell_N}$ $\ell^N$ be \unsure{the space of}}\inserted{For \deleted{every}} \deleted{\unsure{real}} complex sequence $x=(x_i)$ \deleted{such that \unsure{the suit}} \deleted{\unsure{we set|$N\quad\rho$}} \deleted{let $\unsure{N(x\mathord{-}x|N(x\,x}$}

\[
\deleted{N(|x_1|,\mathord{\cdot}\,\mathord{-}|x_n|,0,0,\mathord{-}\,\mathord{-}\,\mathord{\cdot})\quad\text{increases with }n,\quad\text{we set.}}
\]

Let $|x|_n=(|x_1|,\mathord{\cdot}\,\mathord{-}|x_n|,0,\mathord{-}\,\mathord{-}0)$. \deleted{\unsure{We immediately have|Since we have}}\inserted{If $n$ varies} \deleted{then $|x+y|_n\leq |x|_n+|y|_n$} $|x|_n$ is an increasing sequence of positive elements of $R^{(N)}$, let us set

\[
\text{(4)}\qquad N(x)=\lim_n N(|x|_n)
\]

\inserted{this is a number $\geq 0$ finite or infinite} The value \deleted{\unsure{improper|infinite|thus} \unsure{is defined|designated}} of $N(x)$, which coincides with \deleted{\unsure{$N(x)$|that}} \inserted{when $x\in R^{(N)}$.} We have

\[
\text{(5)}\qquad N(\lambda x)=|\lambda|N(x),\qquad N(u+y)\leq N(\unsure{u|x})-N(y)
\]

\deleted{if $N(x)<+\infty$} \inserted{and in} the first formula, \deleted{\unsure{it is to be said|it being said} that \unsure{its second side makes sense}, i.e $|\lambda|$} \inserted{one must take on the second}
\reversemarginpar
\margin{\textcolor{red}{?}}
\normalmarginpar
side the value \textcolor{red}{\uline{\textcolor{black}{zero}}} if it \deleted{\unsure{is in}}\inserted{occurs in the} indeterminate form $0\times\infty$. The second formula follows from the fact that \deleted{$N(x$} $|x+y|_n\leq |x|_n+|y|_n$, whence $N(|x+y|_n)\leq N(|x|_n+|y|_n)\leq N(|x|_n)+N(|y|_n)$.

\deleted{Let} Let $\ell^N$ be the space of complex sequences $x$ such that $N(x)<\infty$, a space \unsure{containing|vector} at least \deleted{\unsure{$C^{(N)}$|those}} the complex sequences zero from \deleted{\unsure{of index|of mult}} a certain index onward. \unsure{On|By|For} $(5)$, we see that $N(x)$ is a \uline{norm} on $\ell^N$.

\deleted{\unsure{Moreover|Besides}, $\ell^N$ is \unsure{a space of} \illegible{2} \unsure{positivity}. Whence \unsure{the injections} $R^{(N)}\subset\ell^N$, and $N(x)$ induces \unsure{by hypothesis}, \unsure{therefore} (even \illegible{1} $\unsure{\alpha^i|\alpha^{(i)}}$), is \illegible{2} \unsure{the candidate|the convergence} in $\ell^N$, \unsure{they therefore necessarily coincide|the convergence at each point}}

\reversemarginpar
\margin{96}
\normalmarginpar
